\section{General weights: what remains after the reductions}

For arbitrary positive $\alpha,\beta,\gamma$ with sum one, the project has not yet found a single identity as rigid as the cubic equal-weight structure.  Nevertheless, the general problem has already been reduced in several useful ways.

\subsection{Interior extremizers have a common Rayleigh normalization}
By the scale-stationary condition, a hypothetical interior extremizer can be normalized so that
\[
B_i=\begin{pmatrix}1&a_i\\ \overline a_i&b_i\end{pmatrix},
\qquad
\alpha a_1+\beta a_2+\gamma a_3=0,
\]
and
\[
\alpha B_1+\beta B_2+\gamma B_3=\diag(1,s).
\]
For the top-eigenvalue problem $s\le1$, whereas for the bottom-eigenvalue problem $s\ge1$.  Thus any genuinely new obstruction must already survive a shared Rayleigh-quotient normalization.

\subsection{Two-dimensional determinant coupling}
The identity
\[
s_1(P)s_2(P)=\det P
\]
means that the $j=1$ and $j=2$ inequalities are not independent in dimension two.  The exact formulation \eqref{eq:hyperbolic-slack} says that both are controlled by a single anisotropy mismatch.  A successful estimate for $L-h$ that uses the determinant slack \eqref{eq:det-slack} would close both indices at once.

\subsection{Rayleigh-envelope candidates}
During the project the following stronger inequalities were identified as natural candidates:
\begin{align}
 s_1(X^\alpha Y^\beta Z^\gamma)
 &\stackrel{?}{\le}
 \max_{\|v\|=1}
 (v^*Xv)^\alpha(v^*Yv)^\beta(v^*Zv)^\gamma,
 \label{eq:Eplus}\\
 s_2(X^\alpha Y^\beta Z^\gamma)
 &\stackrel{?}{\le}
 \min_{\|v\|=1}
 (v^*Xv)^\alpha(v^*Yv)^\beta(v^*Zv)^\gamma.
 \label{eq:Eminus}
\end{align}
If both held, scalar weighted Young applied pointwise in $v$ would prove the Matrix Young inequalities immediately.  However, these envelope statements are stronger than the original target and are currently \statusW{} rather than established lemmas.  They should not be used as hidden assumptions.

\subsection{Rank-one boundary}
\statusW{} The project derived a geometric argument suggesting that three rank-one factors satisfy the $m=3,n=2$ inequality, in contrast with Pang's $m\ge4$ rank-one counterexample mechanism.  The argument passes through the projective geometry of three lines in $\mathbb C^2$.  Because the covering-radius step has not yet received a separate proof audit, this result is recorded as a working lemma rather than as a closed theorem in the present source.

\subsection{Relation to Hua determinant inequalities}
Lin's 2024 paper develops Young-type eigenvalue inequalities alongside extensions of the Hua determinant inequality \cite{Lin2024}.  Its discussion of generalized H\"older eigenvalue inequalities for contractive matrix products is structurally close to the determinant-sign expression
\[
\det(I-P^*P)=1-\Tr(P^*P)+\det(P)^2
\]
that appears in the normalized $2\times2$ $j=2$ problem.  A serious next literature step is therefore to test whether an appropriate contractive rescaling converts the remaining trace-sign inequality into a direct consequence of the Hua/H\"older machinery, rather than expanding matrix cube roots by hand.

\subsection{Current general-weight bottlenecks}
At present the general-weight problem has two main unresolved parts:
\begin{enumerate}[leftmargin=2.2em]
\item prove the top singular-value inequality for genuinely noncommuting, full-rank $2\times2$ triples away from the scalar and commuting boundaries;
\item transfer that control, together with determinant slack, to the bottom singular value, or derive a direct normalized trace-sign inequality analogous to \eqref{eq:J2}.
\end{enumerate}
The equal-weight case should be treated as the model problem from which the correct general-weight scalar inequality is to be inferred.
